Here is a corrected and fully justified proof.

First, one point in the bug report needs to be corrected: the extra assumption \(T\neq\varnothing\) is **necessary**, not optional.  
Indeed, if \(T=\varnothing\), then by definition
\[
\sigma(p)=0\qquad\text{for every }p\in\mathbb R^2,
\]
so
\[
\bigl|\{p\in\mathbb R^2:\sigma(p)\neq 0\}\bigr|=0<n.
\]
Thus part \((1)\) of the statement is false as written. So the theorem that can actually be proved is the same statement with the added hypothesis
\[
T\neq\varnothing.
\]

I will prove that corrected statement.

---

Let
\[
A:=\{p\in\mathbb R^2:\sigma(p)\neq 0\},
\qquad
S+T:=\{s+t:s\in S,\ t\in T\}.
\]
Clearly \(A\subseteq S+T\), so \(A\) is finite.

Call \(p\in S+T\) **lonely** if there is a unique pair \((s,t)\in S\times T\) such that \(p=s+t\).

A basic observation that will be used repeatedly is:

> If \(p=s+t\) is lonely, then the defining sum for \(\sigma(p)\) has exactly one term, so
> \[
> \sigma(p)=\sigma(t)\in\{\pm1\}.
> \]
> In particular, every lonely point belongs to \(A\).

## Case \(n=1\)

Let \(S=\{s\}\).

For each \(t\in T\), the point \(s+t\) has the unique representation \((s,t)\), because the map
\[
t\mapsto s+t
\]
is injective. Hence every point of \(S+T\) is lonely, and
\[
\sigma(s+t)=\sigma(t)\in\{\pm1\}.
\]
Therefore
\[
A=S+T=\{s+t:t\in T\},
\]
so
\[
|A|=|T|\ge 1=n.
\]
This proves part \((1)\) when \(n=1\).

Now assume the extremal case \(|A|=n=1\). Then \(|T|=1\). Write
\[
T=\{t_0\}.
\]
Then:

- for the unique \(s\in S\), the point
  \[
  \ell_s:=s+t_0
  \]
  is lonely, and it is the only lonely point attached to \(s\);
- the set of nonzero points is
  \[
  A=\{\ell_s\};
  \]
- its value is
  \[
  \sigma(\ell_s)=\sigma(t_0)\in\{\pm1\};
  \]
- either \(T_+=\{t_0\}\) and \(T_-=\varnothing\), or \(T_-=\{t_0\}\) and \(T_+=\varnothing\), hence
  \[
  \bigl||T_+|-|T_-|\bigr|=1\le 1.
  \]

So all parts of \((2)\) hold when \(n=1\).

---

## From now on assume \(n\ge 2\)

Fix \(s\in S\), and set
\[
K_s:=\operatorname{conv}(S\setminus\{s\}).
\]

Because \(s\) is a vertex of \(\operatorname{conv}(S)\), it is an extreme point of \(\operatorname{conv}(S)\). Hence
\[
s\notin K_s.
\]
Indeed, if \(s\in K_s\), then \(s\) is a convex combination of points of \(S\setminus\{s\}\). Grouping one term against the convex combination of the remaining terms would write \(s\) as a nontrivial convex combination of two points of \(\operatorname{conv}(S)\), contradicting that \(s\) is an extreme point.

Since \(K_s\) is the convex hull of a finite set, it is compact. Therefore the continuous function
\[
y\mapsto \|s-y\|
\]
attains its minimum on \(K_s\). Choose \(y_s\in K_s\) such that
\[
\|s-y_s\|=\min_{y\in K_s}\|s-y\|.
\]

Now define the linear functional
\[
f_s(x):=(s-y_s)\cdot x.
\]

I claim that
\[
f_s(s)>f_s(x)\qquad\text{for every }x\in K_s.
\]

To prove this, fix \(x\in K_s\). For \(0<\lambda\le 1\), convexity of \(K_s\) gives
\[
y_s+\lambda(x-y_s)\in K_s.
\]
By minimality of \(y_s\),
\[
\|s-y_s\|^2
\le
\bigl\|s-\bigl(y_s+\lambda(x-y_s)\bigr)\bigr\|^2.
\]
Expanding the right-hand side,
\[
\|s-y_s\|^2
\le
\|s-y_s\|^2-2\lambda (s-y_s)\cdot(x-y_s)+\lambda^2\|x-y_s\|^2.
\]
Subtract \(\|s-y_s\|^2\), divide by \(\lambda>0\), and let \(\lambda\downarrow 0\). This yields
\[
(s-y_s)\cdot(x-y_s)\le 0.
\]
Equivalently,
\[
f_s(x)\le f_s(y_s).
\]
On the other hand,
\[
f_s(s)-f_s(y_s)
=
(s-y_s)\cdot(s-y_s)
=
\|s-y_s\|^2
>0,
\]
because \(s\notin K_s\), so \(s\neq y_s\). Therefore
\[
f_s(s)>f_s(y_s)\ge f_s(x)
\qquad\text{for every }x\in K_s.
\]
In particular,
\[
f_s(s)>f_s(u)\qquad\text{for every }u\in S\setminus\{s\}.
\]

Since \(T\) is finite and nonempty, \(f_s\) attains a maximum on \(T\). Choose \(t_s\in T\) such that
\[
f_s(t_s)=\max_{t\in T} f_s(t),
\]
and define
\[
\ell_s:=s+t_s.
\]

### Claim: \(\ell_s\) is lonely

Suppose
\[
\ell_s=u+t
\]
for some \(u\in S\), \(t\in T\). Applying \(f_s\),
\[
f_s(s)+f_s(t_s)=f_s(\ell_s)=f_s(u)+f_s(t).
\]
If \(u\neq s\), then \(u\in S\setminus\{s\}\subseteq K_s\), so \(f_s(u)<f_s(s)\). Hence
\[
f_s(t)=f_s(s)+f_s(t_s)-f_s(u)>f_s(t_s),
\]
contradicting the choice of \(t_s\). Therefore \(u=s\). Then
\[
s+t=s+t_s,
\]
so \(t=t_s\). Thus the representation is unique, and \(\ell_s\) is lonely.

By the observation at the beginning,
\[
\sigma(\ell_s)=\sigma(t_s)\in\{\pm1\},
\]
so \(\ell_s\in A\).

If \(s\neq s'\) and \(\ell_s=\ell_{s'}\), then that common point would have two distinct representations, namely \((s,t_s)\) and \((s',t_{s'})\), contradicting loneliness. Hence the points \(\ell_s\) are pairwise distinct.

Therefore \(A\) contains at least the \(n\) distinct points \(\ell_s\), so
\[
|A|\ge n.
\]
This proves part \((1)\).

---

## Extremal case: assume \(|A|=n\)

We already know that the \(n\) points \(\ell_s\) are distinct and all lie in \(A\). Since \(|A|=n\), it follows that
\[
A=\{\ell_s:s\in S\}.
\]
This proves part \((2)(b)\).

Also, every lonely point belongs to \(A\), so in the extremal case there are no lonely points other than these \(\ell_s\).

### Proof of \((2)(a)\)

Fix \(s\in S\), and suppose \(s+t\) is lonely for some \(t\in T\). Since \(s+t\) is lonely, it lies in \(A\). Hence
\[
s+t=\ell_u
\]
for some \(u\in S\).

If \(u\neq s\), then \(\ell_u\) has two distinct representations:
\[
\ell_u=u+t_u=s+t,
\]
contradicting that \(\ell_u\) is lonely. Therefore \(u=s\). So
\[
s+t=\ell_s=s+t_s,
\]
and hence \(t=t_s\).

Thus for each \(s\in S\), there is exactly one \(t_s\in T\) such that \(s+t_s\) is lonely. This is exactly part \((2)(a)\).

### Proof of \((2)(c)\)

This was already shown:
\[
\sigma(\ell_s)=\sigma(t_s)\in\{\pm1\}
\qquad\text{for every }s\in S.
\]

### Proof of \((2)(d)\)

Because \(\sigma(p)=0\) for \(p\notin S+T\), all relevant sums are finite. We compute
\[
\sum_{p\in S+T}\sigma(p)
=
\sum_{p\in S+T}\ \sum_{\substack{t\in T\\ p-t\in S}}\sigma(t)
=
\sum_{(s,t)\in S\times T}\sigma(t)
=
n\sum_{t\in T}\sigma(t).
\]
Since \(\sigma(t)=+1\) on \(T_+\) and \(\sigma(t)=-1\) on \(T_-\),
\[
\sum_{t\in T}\sigma(t)=|T_+|-|T_-|.
\]
So
\[
\sum_{p\in S+T}\sigma(p)=n\bigl(|T_+|-|T_-|\bigr).
\]

In the extremal case, \(A=\{\ell_s:s\in S\}\), and \(\sigma(p)=0\) for \(p\in (S+T)\setminus A\). Hence
\[
\sum_{p\in S+T}\sigma(p)=\sum_{s\in S}\sigma(\ell_s).
\]
Using \(\sigma(\ell_s)\in\{\pm1\}\),
\[
\left|\sum_{p\in S+T}\sigma(p)\right|
=
\left|\sum_{s\in S}\sigma(\ell_s)\right|
\le
\sum_{s\in S}|\sigma(\ell_s)|
=
n.
\]
Combining the last two displays,
\[
n\,\bigl||T_+|-|T_-|\bigr|\le n.
\]
Since \(n\ge 1\),
\[
\bigl||T_+|-|T_-|\bigr|\le 1.
\]

This proves part \((2)(d)\).

---

So, after the necessary correction \(T\neq\varnothing\), all conclusions follow.  
Without that correction, part \((1)\) is false, so no proof of the original wording can exist.